% Section 9 — Discoveries: the novelty ladder and prospective dossiers (SYNTHETIC results)
\section{What one proposal discovers}
\label{sec:discoveries}

Beating a comparator shows that a first proposal is useful. It does not show that it is new, and
newness is the part of innovation most easily overclaimed. Byte-level novelty is trivial for any
generator; ``absent from the training data'' depends on how hard one looked; and a mechanism is not
discovered because a model emitted it. We therefore grade every committed proposal on a ladder whose
rungs are tested separately (\cref{tab:ladder}) and report every rung with every attempt in the
denominator.

\begin{table}[t]
\caption{The novelty ladder. Each rung adds a test to the one below. Rungs N0 to N4 are computed for
every committed proposal against a complete-artifact archive frozen before unblinding; N5 and N6
are decided only in prospective dossiers.}
\label{tab:ladder}
\centering\tablesf
\begin{tabularx}{\linewidth}{@{}llL@{}}
\toprule
Rung & Name & Test \\
\midrule
N0 & new bytes & the artifact differs from every archived artifact as bytes \\
N1 & structurally new & no archived artifact is equivalent under canonicalization and an enumerated set of rewrite rules, frozen before unblinding \\
N2 & outside the archive & no archived artifact lies within a predeclared behavioural distance on the development problems \\
N3 & useful & verified utility above the locked comparator on the problem for which it was proposed \\
N4 & useful outside the archive & N2 and N3 on the same problem \\
N5 & mechanism supported & on reserved problems, removing the new component and replacing it with the nearest known alternative at matched capacity both lose utility, with interval bounds above zero; prior-work review finds no equivalent \\
N6 & independently replicated & a group not involved in the proposal reproduces the N5 effect with its own evaluation setup, data or fabrication \\
\bottomrule
\end{tabularx}
\end{table}

\subsection{How often a first proposal is new and useful}

Of \syn[EZ-G20]{\EZctxA} committed learning procedures, \syn[EZ-G15]{\EZnovAone\%} are structurally
new, \syn[EZ-G15]{\EZnovAtwo\%} lie outside the archive, and \syn[EZ-G15]{\EZnovAfour\%}
[\syn[EZ-G15]{\EZnovAfourLo}, \syn[EZ-G15]{\EZnovAfourHi}] are both outside it and better than the
comparator (\cref{fig:discovery}). For CAD the rates are \syn[EZ-G15]{\EZnovCone\%},
\syn[EZ-G15]{\EZnovCtwo\%} and \syn[EZ-G15]{\EZnovCfour\%} of \syn[EZ-G20]{\EZspecC} committed parts.
The denominators include failed proposals. Most first proposals that win are therefore not
reproductions of archived solutions, but N4 is a statement about our archive, not about all prior
art; the rungs above it exist because that distinction matters.

\begin{figure}[t]
\centering
% Figure — the novelty ladder as a funnel (A and C) and the two dossier interventions.
% Reads SYNTHETIC data (gaps EZ-G15, EZ-G16, EZ-G17).
\begin{tikzpicture}
\begin{groupplot}[ez,
    group style={group size=2 by 1, horizontal sep=3.0cm},
    height=5.2cm, axis lines*=left, clip=false,
    legend style={font=\sffamily\fontsize{5.6pt}{6.4pt}\selectfont, draw=none, fill=white,
      fill opacity=0.85, text opacity=1, cells={anchor=west}, row sep=-1.5pt}]
  % ---------------- novelty funnel
  \nextgroupplot[width=0.42\linewidth, title={From committed to replicated},
      xbar, bar width=2.6pt, xmode=log, xmin=0.6, xmax=9000, ymin=0.4, ymax=8.6,
      ytick={1,...,8}, yticklabels={N6 replicated, N5 mechanism supported, dossier opportunity,
        N4 useful outside archive, N2 outside archive, N1 structurally new, valid, committed},
      yticklabel style={font=\sffamily\scriptsize\color{ezink}},
      xtick={1,10,100,1000}, xticklabels={1,10,100,1000}, xlabel={first proposals (log scale)},
      legend pos=south east, nodes near coords, nodes near coords style={font=\sffamily\fontsize{5.2pt}{6pt}\selectfont, text=ezgrey},
      point meta=x]
    \addplot[fill=ezaccent!60, draw=ezaccent] table[col sep=comma, x=A, y expr=\thisrow{y}+0.15] {data/synthetic/funnel.csv};
    \addlegendentry{A learning procedures}
    \addplot[fill=ezwarm!60, draw=ezwarm] table[col sep=comma, x=C, y expr=\thisrow{y}-0.15] {data/synthetic/funnel.csv};
    \addlegendentry{C CAD parts}
  % ---------------- dossier interventions
  \nextgroupplot[width=0.33\linewidth, title={Dossier interventions},
      xmin=-0.12, xmax=0.01, ymin=0.5, ymax=4.5, ytick={1,2,3,4},
      yticklabels={D-C1 nearest known, D-C1 remove feature, D-A1 nearest known, D-A1 remove component},
      yticklabel style={font=\sffamily\scriptsize\color{ezink}},
      xtick={-0.1,-0.05,0}, xticklabel style={/pgf/number format/fixed, /pgf/number format/precision=2},
      xlabel={utility change on 60 reserved problems},
      error bars/error bar style={line width=0.5pt}, error bars/error mark options={rotate=90, mark size=1.4pt}]
    \draw[ezgrey, line width=0.5pt, densely dashed] (axis cs:0,0.5) -- (axis cs:0,4.5);
    \addplot[only marks, mark=*, mark size=1.6pt, ezaccent, ezwhere={case}{D-A1}, error bars/.cd, x dir=both, x explicit]
      table[col sep=comma, x=delta, y=y, x error=err] {data/synthetic/dossier.csv};
    \addplot[only marks, mark=triangle*, mark size=1.8pt, ezwarm, ezwhere={case}{D-C1}, error bars/.cd, x dir=both, x explicit]
      table[col sep=comma, x=delta, y=y, x error=err] {data/synthetic/dossier.csv};
\end{groupplot}
\synfigtag
\end{tikzpicture}

\caption{Left: committed first proposals passing each rung of the novelty ladder, on a logarithmic
scale. Right: the two dossiers' interventions on 60 reserved problems each, with 95\% intervals;
removing the proposed component and replacing it with the nearest known alternative both lose
utility. Synthetic (gaps EZ-G15, EZ-G16, EZ-G17).}
\label{fig:discovery}
\end{figure}

\subsection{Discovery dossiers}

From the N4 proposals we drew \syn[EZ-G15]{\EZdossA} learning-procedure and \syn[EZ-G15]{\EZdossC}
CAD dossier opportunities by a rule fixed before unblinding: among proposals whose new component is
isolable, the lowest hash of the committed artifact is examined first. A dossier states the proposed
mechanism in words, before any intervention, and then tests it on problems reserved for that purpose.
\syn[EZ-G15]{\EZnFiveA} learning-procedure and \syn[EZ-G15]{\EZnFiveC} CAD mechanisms reached N5;
\syn[EZ-G15]{\EZnSix} reached N6. We summarize one dossier per game; all are in \cref{app:dossiers}.

\begin{ezdef}{Dossier D-A1 \quad\normalfont\sffamily\color{ezink}noise-gated category shrinkage \synbadge}
\textit{Statement before testing.} The committed procedure encodes high-cardinality categorical
features by their target means, but shrinks each level toward its parent grouping by an amount set
from the out-of-fold residual variance of a preliminary fit, rather than from the level's count
alone; on noisy targets rare levels are shrunk harder. \textit{Removal} of the gate, keeping a
count-only shrinkage of the same capacity, loses \syn[EZ-G16]{\EZdaRemove}
[\syn[EZ-G16]{\EZdaRemoveLo}, \syn[EZ-G16]{\EZdaRemoveHi}] utility on 60 reserved families.
\textit{Replacement} by the nearest known alternative, smoothing with a cross-validated prior
strength, loses \syn[EZ-G16]{\EZdaReplace}. \textit{Prior-work review} found shrinkage estimators
with fixed or cross-validated strength but no noise-gated variant in the reviewed literature; this
supports N5 and is not a claim that the idea has never been written down.
\end{ezdef}

\begin{ezdef}{Dossier D-C1 \quad\normalfont\sffamily\color{ezink}offset web with load-path gusset \synbadge}
\textit{Statement before testing.} For brackets whose bolt pattern is eccentric to the load line, the
committed part offsets its web toward the load and adds a gusset that carries the bolt-preload path
and the bending path in the same material (\cref{fig:cad}, right). \textit{Removal} of the feature,
with all other dimensions re-optimized within the declared language, raises the mass needed to meet
the stiffness requirement by \syn[EZ-G17]{\EZdcMass\%}, a utility loss of \syn[EZ-G17]{\EZdcRemove}. \textit{Replacement} by the nearest
template rib pattern loses \syn[EZ-G17]{\EZdcReplace} utility. \textit{Physical test.} Of
\syn[EZ-G17]{\EZphysN} parts fabricated and load-tested by a laboratory not involved in the design,
all met the load requirement; measured stiffness differed from the prediction by
\syn[EZ-G17]{\EZphysErr\%} on average.
\end{ezdef}

\subsection{The standard we are holding ourselves to}

AlphaGo's move 37 against Lee Sedol mattered because experts first read it as a mistake and then
changed how they played. The analogue for an innovation game is a first proposal that a prospective
expert panel rated unlikely before seeing its outcome, that survives the N5 interventions, and that
practitioners adopt. N6 is the most this paper can test; adoption is outside any paper's control.
We report how far each discovery climbs and do not round a rung up.
